% ======================================================================
% Parte 4: El dilema ético
% ======================================================================

Cada integrante del equipo deberá seleccionar uno de los tres dilemas (sin repetirse entre integrantes). Cada estudiante debe desarrollar su análisis individual (aproximadamente media cuartilla) justificando su postura ética y profesional.

% ----------------------------------------------------------------------
\subsection*{Dilema 1: Seguridad vs. Privacidad}
% ----------------------------------------------------------------------
\textbf{Pregunta detonadora:} \textit{¿Debe Perfectópolis implementar reconocimiento facial en todas las cámaras si esto pudiera ayudar a localizar delincuentes?}\\[2mm]    
No considero correcto que Perfectópolis implemente reconocimiento facial en todas sus cámaras. Aunque la intención es atrapar delincuentes, creo que esta medida puede traer problemas mucho más graves que los que intenta solucionar. Si bien a primera vista suena bien, pues ayudaría a que la policía y el sistema judicial funcionen de forma más rápida y barata y con esto se solucione un problema muy común en países como México, donde los presuntos culpables pasan demasiado tiempo esperando un juicio porque los tribunales están saturados. Esa misma rapidez es su mayor peligro, pues al depender de un sistema automatizado para agilizar la persecución de delitos, corremos el riesgo de que la ciudadanía sea juzgada por una máquina de forma apresurada y sin un análisis humano real, lo que inevitablemente va a causar una saturación de falsos positivos y detenciones injustas.

Además, con este sistema la gente pierde por completo su anonimato en las calles y hasta en espacios privados. Es un golpe directo al derecho a la privacidad que tantos organismos defienden en nuestra realidad. Por ejemplo, normativas muy importantes como el Reglamento General de Protección de Datos (GDPR) \cite{eu_gdpr_2016} y el marco del NIST Privacy \cite{boeckl2020nist} se basan precisamente en esto: exigir el respeto y la protección de los datos personales para evitar que la tecnología se use para vigilar masivamente a los ciudadanos; con ello, y siguiendo el imperativo categórico de Kant, es inaceptable que se sacrifique a toda la ciudadanía para utilizar sus datos biométricos para detener a algunos delincuentes.

Otro grave problema del uso de vigilancia con reconocimiento facial es la posible discriminación que sectores de la población puedan llegar a sufrir, pues el desarrollo de estos sistemas suele estar lleno de sesgos ideológicos que pueden causar falsos positivos. Por ejemplo, según Bajarin T. \cite{bajarin_ibm_facial}, las poblaciones asiático y afroamericanas fueron 100 veces más propensas a ser identificadas como delincuentes que los hombres blancos; de la misma forma, la población nativoamericana tuvo el ratio de falsos positivos más alto de todas las etnias.

Con esto considero que la implementación de sistemas de reconocimiento y vigilancia masiva no es una opción por vulnerar la privacidad de los ciudadanos, se tiende a discriminar por los sesgos de las tecnologías, además de que no se puede asegurar la transparencia de los implementadores acerca del uso de los datos y la implementación de los sistemas. Con ello, propondría que el uso de la vigilancia con reconocimiento de rostro sea una herramienta estrictamente regulada, limitando su uso a buscar criminales muy específicos de los que se tenga constatada su culpabilidad; de esta manera evitamos prejuzgar a nuestra población y llenarnos de resultados sesgados que saturen el sistema de justicia.

En consecuencia, no se tendrá una herramienta omnisciente capaz de encontrar delincuentes prácticamente al instante, pero se mantendrá el derecho humano de la privacidad, se protegerán los datos de los ciudadanos de empresas externas poco transparentes, así como se evitarán más casos de falsos positivos causados por la discriminación algorítmica, brindando confianza en el sistema judicial a los ciudadanos, siempre y cuando este sea sólido independientemente de las tecnologías que utilizan.
% ----------------------------------------------------------------------
\subsection*{Dilema 2: Automatización vs. Responsabilidad Humana}
% ----------------------------------------------------------------------
\textbf{Pregunta detonadora:} \textit{¿Puede el gobierno permitir que una IA determine automáticamente qué solicitudes ciudadanas son importantes?}\\[2mm]
Considero que delegar toda tarea exclusivamente a sistemas automatizados o inteligencia artificial puede afectar en la calidad de vida de los usuarios de los sistemas que delegan estas decisiones; en este caso, los ciudadanos se ven afectados por sesgos algorítmicos que los modelos de inteligencia artificial sufren con frecuencia, pues como lo vimos anteriormente pueden ver ciertas solicitudes totalmente legítimas clasificadas como de baja prioridad prestando atención a problemas que realmente no afectan la calidad de vida de las personas.

Si bien la implementación de estos algoritmos resulta en una reducción de costos operativos para el gobierno de Perfectópolis, así como una ganancia a quienes se encarguen de implementar este sistema, el beneficio es poco cuando la total automatización implica la pérdida de justicia, equidad que un análisis humano pueda dar, así como la poca transparencia gubernamental y tecnológica; de la misma forma, potencialmente se ve afectado el derecho a petición de la ciudadanía al dar menor prioridad a sus peticiones o invisibilizarlas por sesgos tecnológicos. 

De la misma forma que el punto anterior, no podemos asegurar que los sistemas automatizados sean empáticos ni que carezcan de sesgos tecnológicos e ideológicos que invisibilicen las quejas de ciertos sectores de la población.

Por lo tanto, considero que podrían utilizarse sistemas de inteligencia artificial como una herramienta de preclasificación de las denuncias o como un sistema que brinde asesoría para que la tarea pase a un equipo de expertos humanos quienes determinen la importancia de la denuncia de una manera empática y justa para que puedan proceder a atender las denuncias más rápido que lo harían al usar sistemas no automatizados.

Si bien el proceso es ligeramente más lento y requiere de personal humano para revisar las denuncias, se aumenta la probabilidad de que se dé un trato justo y empático a la población de Perfectópolis, así como que se puede ser más transparente con el proceso de decisión que los sistemas automatizados.
% ----------------------------------------------------------------------
\subsection*{Dilema 3: Servicio Gratuito vs. Explotación de Datos}
% ----------------------------------------------------------------------
\textbf{Pregunta detonadora:} \textit{¿Debe Perfectópolis aceptar que una empresa financie PerfectWiFi a cambio de utilizar datos sobre el comportamiento de los usuarios?}\\[2mm]
La utopía de Perfectópolis tiene un problema al intentar seguirla en un país latinoamericano, y muchos otros países fuera del norte global, pues se tiene una falta grave de inclusión digital y este dilema forma parte de ello pues el gobierno puede fomentar en cierta medida la inclusión digital a partir de asegurar que toda la población, o por lo menos los que tengan acceso a los dispositivos necesarios, pueda acceder a sus sistemas en cualquier momento de manera gratuita y todo esto sin requerir de la enorme financiación que requiere el proveer un servicio de conectividad para un país.

Si bien el gobierno ahorra millones en la infraestructura de la red y los sistemas de Perfectópolis, así como otros sistemas privados, son totalmente accesibles para los ciudadanos de manera gratuita; son los mismos ciudadanos los que se ven afectados por ser convertidos en mercancías por la empresa que provee el servicio, pues los datos de consumo, ubicación y hábitos de navegación son sumamente valiosos para las empresas, lo que probablemente brinde un beneficio desproporcionadamente mayor que el gasto de dar la infraestructura de la red.

Otro problema con esto es que no se puede asegurar la transparencia de la empresa conforme a la seguridad de los datos, así como el trato de estos mismos, con lo que puede darse el caso en el que estos datos sean robados o manipulados por la mala práctica que pueda tener la empresa.

Considero que el uso de los datos, con lo valiosos que pueden llegar a ser, es una recompensa desproporcionada con respecto al beneficio que obtienen los ciudadanos. Sin embargo, en este caso se requiere de brindar la inclusión tecnológica de la mayor cantidad de ciudadanos posible, por ello, buscaría que si se van a utilizar los datos de los ciudadanos, que sean anónimos de manera que la empresa solo pueda recibir métricas de estadísticas globales como cuántas personas hay en una ubicación a cierta hora del día, bloqueando el seguimiento total de los ciudadanos.

Derivado de la decisión lo más probable es que la empresa no acepte el trato y toque buscar a otras empresas más flexibles o tener que financiar la infraestructura con los recursos públicos brindando servicios probablemente más deficientes; sin embargo, aseguramos la integridad de los datos de los ciudadanos.
